\subsection{Mathematical formulations and problem classes}
\label{sec:q1-formulations}

\par
\begingroup
\setlength{\parindent}{0pt}
\setlength{\parskip}{0.5em}

The three models minimize procurement cost plus disutility.
For the community model, this is equivalent to maximizing
community welfare.

Members are indexed by $i\in\mathcal N$, and the horizon
is $\mathcal T=\{0,\ldots,23\}$.
Each time interval lasts one hour.
Loads, PV generation and exchanges are measured in kWh
per interval, while battery powers and contractual
capacity limits are measured in kW.
Conversion from power to energy over each one-hour
interval is implicit, and the one-hour factor is
omitted throughout.

For individually metered members in $\mathcal R$ and $\mathcal I$,
total imports and exports are denoted by
$I_{i,t}=q^{\rm imp}_{i,t}$ and
$E_{i,t}=q^{\rm exp}_{i,t}$.
Member PV generation is $P_{i,t}$, distinct from
community imports and exports
$P_t^{\rm imp},P_t^{\rm exp}$.


\subsubsection{Prosumer model}

The reference load and hourly bounds are
\begin{align}
\ell^{\rm ref}_{i,t}
&=L^{\max}\rho_t,
\label{eq:q1-reference-load}\\
\underline\ell_{i,t}
&=\max\{0,\ell^{\rm ref}_{i,t}-d^{\max}\},
\label{eq:q1-load-lower}\\
\overline\ell_{i,t}
&=\min\{L^{\max},
         \ell^{\rm ref}_{i,t}+d^{\max}\},
\label{eq:q1-load-upper}
\end{align}
where $\rho_t\in[0,1]$ is the reference-load
capacity factor, $L^{\max}=3$~kWh/h is the maximum
consumption rate, and $d^{\max}=1$~kWh is the
maximum hourly deviation.

The minimum daily load equals the daily reference
energy:
\begin{equation}
E_i^{\min}
=\sum_{t\in\mathcal T}\ell^{\rm ref}_{i,t}.
\label{eq:q1-daily-minimum-parameter}
\end{equation}

Positive and negative deviations from the reference
load are denoted by $d^+_{i,t}$ and $d^-_{i,t}$.
The hourly linear disutility is
\begin{equation}
D_i^L(d^+_{i,t},d^-_{i,t})
=c_i^L(d^+_{i,t}+d^-_{i,t}),
\qquad c_i^L>0,
\label{eq:q1-hourly-disutility}
\end{equation}
where $c_i^L$ is measured in DKK/kWh.
Total daily disutility is
\begin{equation}
D_i
=\sum_{t\in\mathcal T}
c_i^L(d^+_{i,t}+d^-_{i,t}).
\label{eq:q1-daily-disutility}
\end{equation}

Every member satisfies
\begin{subequations}
\label{eq:q1-shared-member}
\begin{alignat}{2}
&\underline\ell_{i,t}
\le\ell_{i,t}\le\overline\ell_{i,t},
&\qquad&t\in\mathcal T,
\label{eq:q1-load-bounds}\\
&\sum_{t\in\mathcal T}\ell_{i,t}
\ge E_i^{\min},
&&
\label{eq:q1-daily-energy}\\
&\ell_{i,t}-\ell^{\rm ref}_{i,t}
=d^+_{i,t}-d^-_{i,t},
&\qquad&t\in\mathcal T,
\label{eq:q1-deviation-split}\\
&d^+_{i,t},d^-_{i,t}\ge0,
&\qquad&t\in\mathcal T,
\label{eq:q1-deviation-nonnegative}\\
&0\le P_{i,t}\le\overline P_{i,t},
&\qquad&t\in\mathcal T,
\label{eq:q1-pv-bounds}\\
&0\le p^{\rm ch}_{i,t}\le\overline P_i^{\rm ch},
&\qquad&t\in\mathcal T,
\label{eq:q1-charge-bounds}\\
&0\le p^{\rm dis}_{i,t}\le\overline P_i^{\rm dis},
&\qquad&t\in\mathcal T,
\label{eq:q1-discharge-bounds}\\
&e_{i,0}
=E_i^0
+\rho_i^{\rm ch}p^{\rm ch}_{i,0}
-\frac{p^{\rm dis}_{i,0}}{\rho_i^{\rm dis}},
&&
\label{eq:q1-initial-dynamics}\\
&e_{i,t}
=e_{i,t-1}
+\rho_i^{\rm ch}p^{\rm ch}_{i,t}
-\frac{p^{\rm dis}_{i,t}}{\rho_i^{\rm dis}},
&\qquad&t=1,\ldots,23,
\label{eq:q1-storage-dynamics}\\
&0\le e_{i,t}\le\overline E_i,
&\qquad&t\in\mathcal T,
\label{eq:q1-storage-bounds}\\
&e_{i,23}\ge E_i^0.
&&
\label{eq:q1-terminal-condition}
\end{alignat}
\end{subequations}

We impose $e_{i,23}\ge E_i^0$
to prevent the model from reducing the day's cost by ending with less
stored energy than it started with, while allowing the battery to end
the day above its initial energy level.
Here $\overline P_{i,t}$ is available PV energy,
$E_i^0$ is initial stored energy, $\overline E_i$
is battery capacity, $\overline P_i^{\rm ch}$ and
$\overline P_i^{\rm dis}$ are charging and
discharging power limits, and
$\rho_i^{\rm ch},\rho_i^{\rm dis}$ are the
corresponding efficiencies.

Define member net consumption by
\begin{equation}
q_{i,t}
=\ell_{i,t}-P_{i,t}
+p^{\rm ch}_{i,t}-p^{\rm dis}_{i,t}.
\label{eq:q1-net-consumption}
\end{equation}
Positive $q_{i,t}$ denotes an import requirement,
while negative $q_{i,t}$ denotes a net energy surplus.
In $\mathcal R$ and $\mathcal I$, the metering balance gives
$q_{i,t}=I_{i,t}-E_{i,t}$ using the total member import.

Let $\mathcal X_i$ denote the constraints
\eqref{eq:q1-shared-member} together with
\eqref{eq:q1-net-consumption}.


\subsubsection{Trading prices and contract convention}

The effective full-tariff trading prices are
\begin{equation}
a_t=p_t+\tau^{\rm imp},
\qquad
b_t=p_t-\tau^{\rm exp},
\label{eq:q1-full-prices}
\end{equation}
where $p_t$ is the day-ahead energy price and
$\tau^{\rm imp},\tau^{\rm exp}$ are the import
and export grid tariffs.

Under a contract, the effective prices become
\begin{equation}
a_t^{\rm d}
=p_t+(1-\beta^{\rm DSO})\tau^{\rm imp},
\qquad
b_t^{\rm d}
=p_t-(1-\beta^{\rm DSO})\tau^{\rm exp}.
\label{eq:q1-discounted-prices}
\end{equation}
The discount $\beta^{\rm DSO}\in[0,1]$ applies
only to grid tariffs.
Under a contract, imports within the cap are charged at $a_t^{\rm d}$
and exports are credited at $b_t^{\rm d}$.
Imports above the cap are charged at the day-ahead price plus
the penalty, $p_t+\alpha^{\rm DSO}$.

\par
\endgroup

\subsubsection{Reference model $\mathcal R$}

Without a contract, members independently minimize
procurement cost plus disutility. Members are separately metered
and exchange no energy with their neighbours. Exports from one
member do not offset another member's billed imports.
Let
\[
\Gamma^{\mathcal R}
=
\left\{
\ell_{i,t},\,d^+_{i,t},\,d^-_{i,t},\,P_{i,t},\,
p^{\rm ch}_{i,t},\,p^{\rm dis}_{i,t},\,e_{i,t},\,
I_{i,t},\,E_{i,t}
\right\}_{i\in\mathcal N,\ t\in\mathcal T}
\]
denote the collection of decision variables.
The members' independent problems can equivalently be written together as
\begin{subequations}
\label{eq:q1-model-R}
\begin{align}
J^{\mathcal R}
=\min\limits_{\Gamma^{\mathcal R}}\quad&
\sum_{i\in\mathcal N}\sum_{t\in\mathcal T}
\Big[
c_iP_{i,t}
+a_tI_{i,t}-b_tE_{i,t}
+c_i^L(d^+_{i,t}+d^-_{i,t})
\Big],
\label{eq:q1-R-objective}\\
\mathrm{s.t.}\quad&
\mathcal X_i,
&&i\in\mathcal N,\\
&q_{i,t}-I_{i,t}+E_{i,t}=0,
&&i\in\mathcal N,\ t\in\mathcal T,
\label{eq:q1-R-balance}\\
&I_{i,t},E_{i,t}\ge0,
&&i\in\mathcal N,\ t\in\mathcal T.
\end{align}
\end{subequations}

The contractual cap does not constrain dispatch
or generate penalties in $\mathcal R$.
It is only a reference for measuring congestion.
The aggregated gross import across individual
meters is $\sum_i I_{i,t}^{\mathcal R}$.


\subsubsection{Individual-contract model $\mathcal I$}

Each consumer receives a share of the community cap. The discount
applies to imports within its individual cap and to its exports.
Consumers schedule their operations independently.

First solve $\mathcal R$ and fix the individual
capacity allocation:
\begin{align}
n_i^{\mathcal R}
&=
\left[
\sum_{t\in\mathcal T}
\left(
I_{i,t}^{\mathcal R}-E_{i,t}^{\mathcal R}
\right)
\right]_+,
\label{eq:q1-reference-net-import}\\
\kappa_i
&=\frac{n_i^{\mathcal R}}
        {\sum_{j\in\mathcal N}n_j^{\mathcal R}},
\label{eq:q1-cap-share}\\
\overline P_{i,t}^{\rm DSO}
&=\kappa_i\overline P_t^{\rm DSO},
\label{eq:q1-individual-cap}
\end{align}
where $[y]_+=\max\{y,0\}$.

Let $I_{i,t}^{\rm cap}$ be the import within the individual cap and
$x_{i,t}$ the import above it, both in kWh. The meter records the total
$I_{i,t}=I_{i,t}^{\rm cap}+x_{i,t}$.
The collection of decision variables is
\[
\Gamma^{\mathcal I}
=
\left(
\Gamma^{\mathcal R}
\setminus
\{I_{i,t}\}_{i\in\mathcal N,\ t\in\mathcal T}
\right)
\cup
\{I_{i,t}^{\rm cap},\,x_{i,t}\}_{i\in\mathcal N,\ t\in\mathcal T}.
\]
The individual-contract model is
\begin{subequations}
\label{eq:q1-model-I}

\begin{equation}
\begin{aligned}
J^{\mathcal I}=\min\limits_{\Gamma^{\mathcal I}}\quad&
\sum_{i\in\mathcal N}\sum_{t\in\mathcal T}
\Big[
c_iP_{i,t}
+a_t^{\rm d}I_{i,t}^{\rm cap}
-b_t^{\rm d}E_{i,t}\\
&\qquad\quad
+c_i^L(d^+_{i,t}+d^-_{i,t})
+(p_t+\alpha^{\rm DSO})x_{i,t}
\Big],
\end{aligned}
\label{eq:q1-I-objective}
\end{equation}
\begin{align}
\mathrm{s.t.}\quad&
\mathcal X_i,
&&i\in\mathcal N,\\
&q_{i,t}-I_{i,t}^{\rm cap}-x_{i,t}+E_{i,t}=0,
&&i\in\mathcal N,\ t\in\mathcal T,
\label{eq:q1-I-balance}\\
&0\le I_{i,t}^{\rm cap}
\le\kappa_i\overline P_t^{\rm DSO},
&&i\in\mathcal N,\ t\in\mathcal T,
\label{eq:q1-I-cap}\\
&E_{i,t},x_{i,t}\ge0,
&&i\in\mathcal N,\ t\in\mathcal T.
\end{align}
\end{subequations}

Since $p_t+\alpha^{\rm DSO}>a_t^{\rm d}$ in the base case,
optimality implies
\begin{equation}
x_{i,t}
=
\left[
I_{i,t}
-\kappa_i\overline P_t^{\rm DSO}
\right]_+.
\end{equation}

\subsubsection{Centralized community model $\mathcal C$}

The community manager jointly dispatches all members under a single
contract at one connection point to maximize community welfare.
The community trades with the external grid through a single meter,
while members exchange electricity freely behind the meter.
Only transactions crossing this meter incur external energy trading
costs, grid tariffs and cap-exceedance penalties.

Let $P_t^{\rm imp,cap}$ be the import within the community cap and
$x_t$ the import above it, both in kWh. The common meter records the total
$P_t^{\rm imp}=P_t^{\rm imp,cap}+x_t$.
The collection of decision variables is
\[
\Gamma^{\mathcal C}
=
\left(
\Gamma^{\mathcal I}
\setminus
\{I_{i,t}^{\rm cap},\,x_{i,t},\,E_{i,t}\}_{i\in\mathcal N,\ t\in\mathcal T}
\right)
\cup
\{P_t^{\rm imp,cap},\,x_t,\,P_t^{\rm exp}\}_{t\in\mathcal T}.
\]
The centralized model is
\begin{subequations}
\label{eq:q1-model-C}
\begin{align}
J^{\mathcal C}=\min\limits_{\Gamma^{\mathcal C}}\quad&
\sum_{i\in\mathcal N}\sum_{t\in\mathcal T}
\Big[
c_iP_{i,t}
+c_i^L(d^+_{i,t}+d^-_{i,t})
\Big]
\nonumber\\
&+\sum_{t\in\mathcal T}
\Big[
a_t^{\rm d}P_t^{\rm imp,cap}
-b_t^{\rm d}P_t^{\rm exp}
+(p_t+\alpha^{\rm DSO})x_t
\Big],
\label{eq:q1-C-objective}\\
\mathrm{s.t.}\quad&
\mathcal X_i,
&&i\in\mathcal N,\\
&\sum_{i\in\mathcal N}q_{i,t}
-P_t^{\rm imp,cap}-x_t+P_t^{\rm exp}=0,
&&t\in\mathcal T,
\label{eq:q1-community-balance}\\
&0\le P_t^{\rm imp,cap}
\le\overline P_t^{\rm DSO},
&&t\in\mathcal T,
\label{eq:q1-community-cap}\\
&P_t^{\rm exp},x_t\ge0,
&&t\in\mathcal T.
\end{align}
\end{subequations}

Since $p_t+\alpha^{\rm DSO}>a_t^{\rm d}$ in the base case,
the optimal split satisfies
$x_t=[P_t^{\rm imp}-\overline P_t^{\rm DSO}]_+$.

\subsubsection{Problem classes and consequences for duality}

The objective functions and all constraints of the three
models are linear, and there are no integer variables.
In the base case, each model is feasible and its objective
is bounded below. Therefore, strong duality holds:
\begin{equation}
J^{\mathcal R}=D^{\mathcal R},\qquad
J^{\mathcal I}=D^{\mathcal I},\qquad
J^{\mathcal C}=D^{\mathcal C}.
\end{equation}
The KKT conditions are necessary and sufficient for
global optimality.
